\documentclass[a4paper,10pt,english]{article}

\usepackage[margin=2.5cm]{geometry}
\usepackage[utf8]{inputenc}
\usepackage[T1]{fontenc}
\usepackage[english]{babel}
\usepackage{amsmath,amssymb,amsthm}
\usepackage[numbers,sort&compress]{natbib}
\usepackage[colorlinks,urlcolor=blue,citecolor=blue,linkcolor=blue]{hyperref}
\usepackage{review-mode}
\reviewon
\commenton
\todoon

\theoremstyle{definition}
\newtheorem{assumption}{Assumption}
\newtheorem{lemma}{Lemma}
\newtheorem{proposition}{Proposition}
\newtheorem{theorem}[proposition]{Theorem}
% Continue the principal-result numbering after the first lemma.
\setcounter{proposition}{1}
\DeclareMathOperator{\Var}{Var}
\DeclareMathOperator{\Cov}{Cov}
\DeclareMathOperator{\rank}{rank}
\setlength{\emergencystretch}{3em}
\renewcommand{\ReviewRevision}{appendix-new-2026-10-31}

\begin{document}
\begin{center}
\Large Appendix: Uniform Moment and Denominator Bounds
\end{center}
\medskip

\begin{review}
We observe
\begin{equation}\label{eq:new-model}
 y^{(i)}
 =\beta_0x_0^{(i)}
 +{X_{-0}^{(i)}}^\top\beta_{-0}
 +\varepsilon^{(i)},
 \qquad i=1,\ldots,n.
\end{equation}
\rbreak
Write
\[
 \sigma_0^2=\Var(x_0),\qquad
 C=\Cov(X_{-0}),\qquad
 d=\Cov(X_{-0},x_0),\qquad
 \sigma_\varepsilon^2=\Var(\varepsilon).
\]
\rbreak
Expectations refer to one population observation, and
\(\|V\|_{\psi_2}
 =\inf\{a>0:\mathbb E[\exp(V^2/a^2)]\le2\}\).

\begin{assumption}\label{ass:new}
The following conditions hold.
\begin{enumerate}
\renewcommand{\labelenumi}{}
\item \theenumi.\ Independent observations:
\begin{equation*}\tag{A1}\label{ass:new-iid}
 (x_0^{(i)},X_{-0}^{(i)},\varepsilon^{(i)}),
 \qquad i=1,\ldots,n,
 \quad\text{are iid.}
\end{equation*}
\item \theenumi.\ Subgaussian design:
\begin{equation*}\tag{A2}\label{ass:new-design}
 \mathbb E[(x_0,X_{-0})]=0,\qquad
 \sigma_0^2>0,\qquad
 \sup_{\substack{u\in\mathbb R^{p+1}\\
 \mathbb E[\langle(x_0,X_{-0}),u\rangle^2]=1}}
 \|\langle(x_0,X_{-0}),u\rangle\|_{\psi_2}
 \le K_X.
\end{equation*}
\item \theenumi.\ Independent noise:
\begin{equation*}\tag{A3}\label{ass:new-noise}
 \begin{gathered}
 \varepsilon\perp(x_0,X_{-0}),\qquad
 \mathbb E[\varepsilon]=0,\qquad
 \|\varepsilon\|_{\psi_2}\le K_\varepsilon\sigma_\varepsilon,\\
 \beta_{-0}^\top C\beta_{-0}+\sigma_\varepsilon^2>0.
 \end{gathered}
\end{equation*}
\item \theenumi.\ Local overlap:
\begin{equation*}\tag{A4}\label{ass:new-overlap}
 \min_{\substack{S\subseteq[p]\\|S|\le s}}
 \left(\sigma_0^2-d_S^\top C_{SS}^{-1}d_S\right)
 \ge c_{\mathrm{ov}}\sigma_0^2.
\end{equation*}
\end{enumerate}
The constants \(K_X,K_\varepsilon\ge1\) and
\(c_{\mathrm{ov}}>0\) are fixed, and \(1\le s\le p\).
\rbreak
Inverses of singular covariance matrices mean Moore--Penrose inverses.
\end{assumption}

For \(S\subseteq[p]\), the OLS coefficients satisfy
\begin{equation}\label{eq:new-normal}
 \begin{pmatrix}
 x_0^\top x_0&x_0^\top X_S\\
 X_S^\top x_0&X_S^\top X_S
 \end{pmatrix}
 \begin{pmatrix}\hat\beta_{0,S}\\\hat\gamma_S\end{pmatrix}
 =
 \begin{pmatrix}x_0^\top y\\X_S^\top y\end{pmatrix}.
\end{equation}
\rbreak
For \(x_0^\top M_Sx_0>0\),
\begin{equation}\label{eq:new-coefficient}
 \hat\beta_{0,S}
 =\frac{x_0^\top M_Sy}{x_0^\top M_Sx_0},
 \qquad M_S=I-P_S,
 \qquad P_S=X_S(X_S^\top X_S)^\dagger X_S^\top.
\end{equation}
\rbreak
On the event \(x_0^\top M_Sx_0=0\), the OLS treatment coefficient
is not uniquely identified.

To bound the numerator and denominator in
\eqref{eq:new-coefficient}, define
\begin{equation}\label{eq:new-augmented}
 \begin{aligned}
 \widehat\Sigma_J
 &=\frac1n
 \begin{pmatrix}x_0&X_J&y\end{pmatrix}^{\!\top}
 \begin{pmatrix}x_0&X_J&y\end{pmatrix},\\
 \Sigma_J
 &=\mathbb E\left[
 \begin{pmatrix}x_0\\X_J\\y\end{pmatrix}
 \begin{pmatrix}x_0\\X_J\\y\end{pmatrix}^{\!\top}
 \right],
 \end{aligned}
\end{equation}
and, for \(k=\min(2s,p)\) and \(t\ge1\), set
\begin{equation}\label{eq:new-radius}
 r_t=C_0(K_X^2+K_\varepsilon^2)
 \left[
 \sqrt{\frac{k+2}{n}}
 +\sqrt{\frac{\log\binom pk+t}{n}}
 \right],
\end{equation}
where \(C_0\) is a universal constant, and let
\begin{equation}\label{eq:new-event}
 \mathcal E_t
 =\bigcap_{\substack{J\subseteq[p]\\|J|\le k}}
 \left\{
 (1-r_t)\Sigma_J\preceq\widehat\Sigma_J
 \preceq(1+r_t)\Sigma_J
 \right\},
\end{equation}
where the order \(A\preceq B\) means
\(u^\top Au\le u^\top Bu\) for every vector \(u\).
\rbreak
The leading block and last column of \(\widehat\Sigma_J\) contain,
respectively, the matrix and right-hand side of \eqref{eq:new-normal}.
\rbreak
The augmented matrix \([x_0,X_{-0},y]\) has \(p+2\) columns,
of which \(|J|+2\) remain after selecting the controls in \(J\).

\begin{lemma}[Uniform moment and denominator bounds]\label{lem:new-uniform}
There are universal constants \(C_0,C_1\) such that, under
Assumption~\ref{ass:new} and
\begin{equation}\label{eq:new-sample-size}
 n\ge C_1(K_X^2+K_\varepsilon^2)^2
 \left[k+2+\log\binom pk+t\right],
\end{equation}
with probability at least \(1-2e^{-t}\), the event \(\mathcal E_t\)
holds and, simultaneously for every \(|J|\le s\),
\begin{equation}\label{eq:new-denominator}
 \begin{aligned}
 &(1-r_t)\left(\sigma_0^2-d_J^\top C_{JJ}^{-1}d_J\right)\\
 &\qquad\le\frac{x_0^\top M_Jx_0}{n}
 \le(1+r_t)\left(\sigma_0^2-d_J^\top C_{JJ}^{-1}d_J\right).
 \end{aligned}
\end{equation}
The radius satisfies \(r_t\le1/2\), so every searched coefficient
is uniquely identified on this event.
\end{lemma}

\begin{proof}
Fix \(J\).
\rbreak
Adding and subtracting \(c\beta_0x_0\), and using
\eqref{eq:new-model}, gives
\[
 \begin{aligned}
 a x_0+X_J^\top v+cy
 &=(a+c\beta_0)x_0+X_J^\top v+c(y-\beta_0x_0)\\
 &=(a+c\beta_0)x_0+X_J^\top v
 +cX_{-0}^\top\beta_{-0}+c\varepsilon.
 \end{aligned}
\]
\rbreak
With \(K_*=(K_X^2+K_\varepsilon^2)^{1/2}\),
\[
 \begin{aligned}
 &\|a x_0+X_J^\top v+cy\|_{\psi_2}\\
 &\le
 \|(a+c\beta_0)x_0+X_J^\top v
 +cX_{-0}^\top\beta_{-0}\|_{\psi_2}
 +|c|\|\varepsilon\|_{\psi_2}\\
 &\le
 K_X\|(a+c\beta_0)x_0+X_J^\top v
 +cX_{-0}^\top\beta_{-0}\|_{L_2}
 +K_\varepsilon|c|\sigma_\varepsilon\\
 &\le
 K_*\left(
 \mathbb E\left[
 \bigl((a+c\beta_0)x_0+X_J^\top v
 +cX_{-0}^\top\beta_{-0}\bigr)^2
 \right]+c^2\sigma_\varepsilon^2
 \right)^{1/2}\\
 &=K_*\|a x_0+X_J^\top v+cy\|_{L_2}.
 \end{aligned}
\]
The first inequality is the triangle inequality for the subgaussian norm;
the second uses \eqref{ass:new-design}--\eqref{ass:new-noise};
the third is Cauchy--Schwarz in \(\mathbb R^2\);
and the final equality uses
\(\varepsilon\perp(x_0,X_{-0})\) and
\(\mathbb E[\varepsilon]=0\).
\rbreak
We have therefore shown that the augmented row
\((x_0,X_{-0}^\top,y)^\top\), and each selected row
\((x_0,X_J^\top,y)^\top\), are subgaussian relative to their
population covariances, with constant at most \(K_*\).
\rbreak
This is the bound needed to apply a covariance estimate to
\(\widehat\Sigma_J\), whose population counterpart is \(\Sigma_J\).

We now apply \citet[Theorem~4.6.1 and equation~(4.22)]{vershynin2018hdp},
which gives
\begin{equation}\label{eq:new-fixed-support}
 \mathbb P\left(
 \begin{gathered}
 (1-\rho)\Sigma_J\preceq\widehat\Sigma_J
 \preceq(1+\rho)\Sigma_J\\[-2pt]
 \text{fails}
 \end{gathered}
 \right)\le2e^{-u},
\end{equation}
where
\(\rho=CK_*^2[\sqrt{(|J|+2+u)/n}+(|J|+2+u)/n]\).
\rbreak
In that result, \(m\) is our sample size \(n\), the column dimension
\(n\) is \(\rank(\Sigma_J)\le |J|+2\), \(A\) is the
population-normalized augmented matrix \([x_0,X_J,y]\) on its covariance
range, \(K\le K_*\), and the tail parameter \(t\) is our \(\sqrt u\).
\rbreak
Choose \(C_0\ge\max\{1,2C\}\) and \(C_1=8C_0^2\), where \(C\)
is the universal constant in the displayed expression for \(\rho\).
\rbreak
The reference leaves the numerical value of \(C\) unspecified;
all dependence on the subgaussian norms is explicit in
\eqref{eq:new-radius} and \eqref{eq:new-sample-size}.

It suffices to consider sets with \(|J|=k\), because every smaller
augmented matrix is a principal submatrix of one of these.
\rbreak
Choose \(u=t+\log\binom pk\).
\rbreak
Using \eqref{eq:new-sample-size}, the threshold on the right-hand side
of the comparison in \eqref{eq:new-fixed-support} satisfies
\begin{equation}\label{eq:new-threshold}
 \begin{aligned}
 \rho
 &\le CK_*^2\left[
 \sqrt{\frac{k+2+\log\binom pk+t}{n}}
 +\frac{k+2+\log\binom pk+t}{n}
 \right]\\
 &\le2CK_*^2\sqrt{\frac{k+2+\log\binom pk+t}{n}}\\
 &\le C_0K_*^2\left[
 \sqrt{\frac{k+2}{n}}
 +\sqrt{\frac{\log\binom pk+t}{n}}
 \right],
 \end{aligned}
\end{equation}
where the right-hand side of \eqref{eq:new-threshold} is \(r_t\)
by \eqref{eq:new-radius}.
\rbreak
A union bound gives
\begin{equation}\label{eq:new-probability}
 \mathbb P(\mathcal E_t^c)
 \le\sum_{|J|=k}2e^{-t-\log\binom pk}
 =2e^{-t},
\end{equation}
which proves the probability bound in the lemma.

Focus now on the value of \(r_t\).
\rbreak
Using \((a+b)^2\le2(a^2+b^2)\) in \eqref{eq:new-radius}, followed by
\eqref{eq:new-sample-size} and \(C_1=8C_0^2\), gives
\begin{equation}\label{eq:new-radius-small}
 \begin{aligned}
 r_t^2
 &\le\frac{2C_0^2K_*^4}{n}
 \left[k+2+\log\binom pk+t\right]\\
 &\le\frac{2C_0^2}{C_1}=\frac14.
 \end{aligned}
\end{equation}

Assume now that we are on the event \(\mathcal E_t\), and fix
\(|J|\le s\) and \(a\in\mathbb R^{|J|}\).
\rbreak
The coefficient vector \((1,-a^\top,0)^\top\) gives
\begin{equation}\label{eq:new-residual-comparison}
 \begin{aligned}
 &(1-r_t)\mathbb E[(x_0-X_J^\top a)^2]\\
 &\qquad\le\frac{\|x_0-X_Ja\|^2}{n}
 \le(1+r_t)\mathbb E[(x_0-X_J^\top a)^2].
 \end{aligned}
\end{equation}
\rbreak
We now use
\[
 \min_a\frac{\|x_0-X_Ja\|^2}{n}
 =\frac{x_0^\top M_Jx_0}{n},
\]
by the definition of the orthogonal projection onto the column span of
\(X_J\).
\rbreak
The population minimum is attained at \(a_*=C_{JJ}^{-1}d_J\), since
\[
 \mathbb E[(x_0-X_J^\top a)^2]
 =\sigma_0^2-d_J^\top C_{JJ}^{-1}d_J
 +(a-a_*)^\top C_{JJ}(a-a_*).
\]
\rbreak
Minimizing the lower inequality in
\eqref{eq:new-residual-comparison}, and evaluating its upper inequality
at \(a_*\), proves the two-sided bound \eqref{eq:new-denominator}.
\rbreak
Applying \eqref{ass:new-overlap} to its lower bound gives
\[
 \inf_{|J|\le s}\frac{x_0^\top M_Jx_0}{n}
 \ge(1-r_t)c_{\mathrm{ov}}\sigma_0^2>0.\qedhere
\]
\end{proof}

For the subsequent coefficient bounds, define
\begin{equation}\label{eq:new-delta}
 \delta_{s,n,p}=\sqrt{\frac{s\log(ep/s)}n}.
\end{equation}
\rbreak
Some elementary algebra
(\(k+2+\log\binom pk\le6s\log(ep/s)\)) then yields
\begin{equation}\label{eq:new-radius-rate}
 r_t\le C_0(K_X^2+K_\varepsilon^2)
 \left(\sqrt{12}\,\delta_{s,n,p}+\sqrt{\frac{2t}{n}}\right).
\end{equation}

Write
\[
 w=y-\beta_0x_0,\qquad
 \sigma_w^2=\mathbb E[w^2]
 =\beta_{-0}^\top C\beta_{-0}+\sigma_\varepsilon^2,\qquad
 \eta(S)=\sigma_0^2-d_S^\top C_{SS}^{-1}d_S.
\]
\rbreak
The quantity \(\eta(S)\) is the population residual variance of the
treatment after linear adjustment on \(X_S\).
\rbreak
Assumption~\eqref{ass:new-overlap} requires
\(\eta(S)\ge c_{\mathrm{ov}}\sigma_0^2\) for every \(|S|\le s\).

The coefficient of \(x_0\) in the population linear projection of \(y\)
on \((x_0,X_S)\) is
\begin{equation}\label{eq:new-population-coefficient}
 \beta^\star_{0,S}=\beta_0+b(S),\qquad
 b(S)=
 \frac{
 d^\top\beta_{-0}
 -d_S^\top C_{SS}^{-1}(C\beta_{-0})_S
 }{\eta(S)},
\end{equation}
with \(b(\emptyset)=d^\top\beta_{-0}/\sigma_0^2\).
\rbreak
Consequently,
\begin{equation}\label{eq:new-population-change}
 \eta(S)\bigl(b(S)-b(\emptyset)\bigr)
 =
 -d_S^\top C_{SS}^{-1}(C\beta_{-0})_S
 +b(\emptyset)d_S^\top C_{SS}^{-1}d_S.
\end{equation}

\begin{proposition}\label{prop:new-coefficient}
Under Assumption~\ref{ass:new} and the sample-size condition
\eqref{eq:new-sample-size}, on the event \(\mathcal E_t\), simultaneously
for every \(S\subseteq[p]\) with \(|S|\le s\),
\begin{equation}\label{eq:new-coefficient-bound}
 \begin{aligned}
 &\left|
 (\hat\beta_{0,S}-\hat\beta_{0,\emptyset})
 -\bigl(b(S)-b(\emptyset)\bigr)
 \right|\\
 &\quad\le
 \frac{A r_t}{\eta(S)}
 \left[
 \sigma_0
 \sqrt{(C\beta_{-0})_S^\top C_{SS}^{-1}(C\beta_{-0})_S}
 +\sigma_w\sqrt{d_S^\top C_{SS}^{-1}d_S}
 +r_t\sigma_0\sigma_w
 \right],
 \end{aligned}
\end{equation}
where \(A\) is a universal numerical constant.
\end{proposition}

\begin{proof}
Fix \(S\) and work on \(\mathcal E_t\).
\rbreak
By the definition of \(w\) and the OLS formula
\eqref{eq:new-coefficient},
\begin{equation}\label{eq:new-sample-change}
 \begin{aligned}
 \hat\beta_{0,S}-\hat\beta_{0,\emptyset}
 &=(\hat\beta_{0,S}-\beta_0)
 -(\hat\beta_{0,\emptyset}-\beta_0)\\
 &=
 \frac{x_0^\top M_Sw}{x_0^\top M_Sx_0}
 -(\hat\beta_{0,\emptyset}-\beta_0)\\
 &=
 \frac{
 x_0^\top w-x_0^\top P_Sw
 -(\hat\beta_{0,\emptyset}-\beta_0)
 (x_0^\top x_0-x_0^\top P_Sx_0)
 }{x_0^\top M_Sx_0}\\
 &=
 \frac{
 -x_0^\top P_Sw
 +(\hat\beta_{0,\emptyset}-\beta_0)x_0^\top P_Sx_0
 }{x_0^\top M_Sx_0}.
 \end{aligned}
\end{equation}
The last equality uses the normal equation
\(x_0^\top w=(\hat\beta_{0,\emptyset}-\beta_0)x_0^\top x_0\).

Multiply \eqref{eq:new-sample-change} by \(x_0^\top M_Sx_0/n\),
subtract \eqref{eq:new-population-change}, and add and subtract
\((\hat\beta_{0,\emptyset}-\beta_0)d_S^\top C_{SS}^{-1}d_S\).
\rbreak
This yields
\begin{equation}\label{eq:new-four-terms}
 \begin{aligned}
 &(\hat\beta_{0,S}-\hat\beta_{0,\emptyset})
 -\bigl(b(S)-b(\emptyset)\bigr)\\
 &=\frac{1}{x_0^\top M_Sx_0/n}
 \Bigg[
 \underbrace{
 -\left\{\frac{x_0^\top P_Sw}{n}
 -d_S^\top C_{SS}^{-1}(C\beta_{-0})_S\right\}
 }_{(T_1)}\\
 &\qquad+
 \underbrace{
 (\hat\beta_{0,\emptyset}-\beta_0)
 \left\{\frac{x_0^\top P_Sx_0}{n}
 -d_S^\top C_{SS}^{-1}d_S\right\}
 }_{(T_2)}\\
 &\qquad+
 \underbrace{
 (\hat\beta_{0,\emptyset}-\beta_0-b(\emptyset))
 d_S^\top C_{SS}^{-1}d_S
 }_{(T_3)}\\
 &\qquad-
 \underbrace{
 \bigl(b(S)-b(\emptyset)\bigr)
 \left\{\frac{x_0^\top M_Sx_0}{n}-\eta(S)\right\}
 }_{(T_4)}
 \Bigg].
 \end{aligned}
\end{equation}

We now control each of the \(T\) terms and the denominator
\(x_0^\top M_Sx_0/n\).
\rbreak
Lemma~\ref{lem:new-uniform} gives
\[
 \frac{x_0^\top M_Sx_0}{n}
 \ge(1-r_t)\eta(S)\ge\frac{\eta(S)}2,\qquad
 \left|\frac{x_0^\top M_Sx_0}{n}-\eta(S)\right|
 \le r_t\eta(S).
\]
\rbreak
Using \eqref{eq:new-population-change},
\begin{equation}\label{eq:new-fourth-term}
 \begin{aligned}
 |T_4|
 &\le r_t\eta(S)|b(S)-b(\emptyset)|\\
 &=r_t\left|
 d_S^\top C_{SS}^{-1}
 \bigl((C\beta_{-0})_S-b(\emptyset)d_S\bigr)
 \right|\\
 &\le r_t\left[
 \sqrt{d_S^\top C_{SS}^{-1}d_S}\,
 \sqrt{(C\beta_{-0})_S^\top C_{SS}^{-1}(C\beta_{-0})_S}
 +|b(\emptyset)|\,d_S^\top C_{SS}^{-1}d_S
 \right]\\
 &\le r_t\left[
 \sigma_0
 \sqrt{(C\beta_{-0})_S^\top C_{SS}^{-1}(C\beta_{-0})_S}
 +\sigma_w\sqrt{d_S^\top C_{SS}^{-1}d_S}
 \right],
 \end{aligned}
\end{equation}
where the penultimate inequality is Cauchy--Schwarz and the triangle
inequality, and the last uses
\(\sqrt{d_S^\top C_{SS}^{-1}d_S}\le\sigma_0\) and
\(|b(\emptyset)|\le\sigma_w/\sigma_0\).

For \(T_1\) and the projected moment
\((x_0^\top P_Sx_0/n)\) in \(T_2\), apply
Lemma~\ref{lem:new-projected}(b) with \(z=w\) and \(z=x_0\), respectively.
\rbreak
Apply Lemma~\ref{lem:new-projected}(a) to the coefficient factors in
\(T_2\) and \(T_3\).
\rbreak
Thus
\begin{equation}\label{eq:new-assembly}
 \begin{aligned}
 &\left|
 (\hat\beta_{0,S}-\hat\beta_{0,\emptyset})
 -\bigl(b(S)-b(\emptyset)\bigr)
 \right|\\
 &\le
 \frac{2}{\eta(S)}
 \bigl(|T_1|+|T_2|+|T_3|+|T_4|\bigr)\\
 &\le
 \frac{2r_t}{\eta(S)}
 \Bigg[
 4\sigma_0
 \sqrt{(C\beta_{-0})_S^\top C_{SS}^{-1}(C\beta_{-0})_S}
 +4\sigma_w\sqrt{d_S^\top C_{SS}^{-1}d_S}
 +4r_t\sigma_0\sigma_w\\
 &\hspace{22mm}
 +12\sigma_w
 \left(2\sqrt{d_S^\top C_{SS}^{-1}d_S}+r_t\sigma_0\right)
 +4\frac{\sigma_w}{\sigma_0}d_S^\top C_{SS}^{-1}d_S\\
 &\hspace{22mm}
 +\sigma_0
 \sqrt{(C\beta_{-0})_S^\top C_{SS}^{-1}(C\beta_{-0})_S}
 +\sigma_w\sqrt{d_S^\top C_{SS}^{-1}d_S}
 \Bigg]\\
 &\le
 \frac{A r_t}{\eta(S)}
 \left[
 \sigma_0
 \sqrt{(C\beta_{-0})_S^\top C_{SS}^{-1}(C\beta_{-0})_S}
 +\sigma_w\sqrt{d_S^\top C_{SS}^{-1}d_S}
 +r_t\sigma_0\sigma_w
 \right].
 \end{aligned}
\end{equation}
The last inequality uses
\(d_S^\top C_{SS}^{-1}d_S
\le\sigma_0\sqrt{d_S^\top C_{SS}^{-1}d_S}\).
\end{proof}

\begin{lemma}\label{lem:new-projected}
On \(\mathcal E_t\), with \(r_t\le1/2\), the following bounds hold
simultaneously for every \(|S|\le s\).

\textup{(a)} The unadjusted coefficient satisfies
\begin{equation}\label{eq:new-baseline-bounds}
 \begin{aligned}
 |\hat\beta_{0,\emptyset}-\beta_0-b(\emptyset)|
 &\le4r_t\frac{\sigma_w}{\sigma_0},\\
 |\hat\beta_{0,\emptyset}-\beta_0|
 &\le3\frac{\sigma_w}{\sigma_0}.
 \end{aligned}
\end{equation}

\textup{(b)} For \(z=x_0\) and \(z=w\),
\begin{equation}\label{eq:new-projected-bound}
 \begin{aligned}
 &\left|
 \frac{x_0^\top P_Sz}{n}
 -d_S^\top C_{SS}^{-1}\mathbb E[X_Sz]
 \right|\\
 &\quad\le
 4r_t\left[
 \sigma_0
 \sqrt{\mathbb E[X_Sz]^\top C_{SS}^{-1}\mathbb E[X_Sz]}
 +\|z\|_{L_2}\sqrt{d_S^\top C_{SS}^{-1}d_S}
 +r_t\sigma_0\|z\|_{L_2}
 \right].
 \end{aligned}
\end{equation}
\end{lemma}

\begin{proof}
By Lemma~\ref{lem:new-uniform}, on \(\mathcal E_t\),
\begin{equation}\label{eq:new-augmented-operator}
 \left\|
 \Sigma_S^{-1/2}
 (\widehat\Sigma_S-\Sigma_S)
 \Sigma_S^{-1/2}
 \right\|_{\mathrm{op}}
 \le r_t.
\end{equation}
\rbreak
For every deterministic pair \(u,v\in\mathbb R^{|S|+2}\),
Cauchy--Schwarz therefore gives
\begin{equation}\label{eq:new-bilinear}
 \begin{aligned}
 \left|u^\top(\widehat\Sigma_S-\Sigma_S)v\right|
 &=
 \left|
 (\Sigma_S^{1/2}u)^\top
 \left[
 \Sigma_S^{-1/2}(\widehat\Sigma_S-\Sigma_S)\Sigma_S^{-1/2}
 \right]
 (\Sigma_S^{1/2}v)
 \right|\\
 &\le
 r_t\sqrt{u^\top\Sigma_Su}\,
 \sqrt{v^\top\Sigma_Sv}.
 \end{aligned}
\end{equation}

For part (a), take \(u=(1,0_{|S|}^\top,0)^\top\) and
\(v=(-\beta_0,0_{|S|}^\top,1)^\top\).
\rbreak
Then \([x_0,X_S,y]u=x_0\), \([x_0,X_S,y]v=w\),
\(u^\top\Sigma_Su=\sigma_0^2\), and
\(v^\top\Sigma_Sv=\sigma_w^2\).
\rbreak
Applying \eqref{eq:new-bilinear} to \(u,v\) and to \(u,u\) gives
\begin{equation}\label{eq:new-scalar-moments}
 \begin{aligned}
 \left|\frac{x_0^\top w}{n}-\mathbb E[x_0w]\right|
 &=
 \left|u^\top(\widehat\Sigma_S-\Sigma_S)v\right|
 \le r_t\sigma_0\sigma_w,\\
 \left|\frac{x_0^\top x_0}{n}-\sigma_0^2\right|
 &=
 \left|u^\top(\widehat\Sigma_S-\Sigma_S)u\right|
 \le r_t\sigma_0^2.
 \end{aligned}
\end{equation}

By the triangle inequality,
\[
 \begin{aligned}
 &|\hat\beta_{0,\emptyset}-\beta_0-b(\emptyset)|\\
 &=
 \frac{
 |x_0^\top w/n-b(\emptyset)x_0^\top x_0/n|
 }{x_0^\top x_0/n}\\
 &\le
 \frac{
 |x_0^\top w/n-\mathbb E[x_0w]|
 +|b(\emptyset)|\,|x_0^\top x_0/n-\sigma_0^2|
 }{x_0^\top x_0/n}\\
 &\le
 \frac{
 r_t\sigma_0\sigma_w+r_t|b(\emptyset)|\sigma_0^2
 }{(1-r_t)\sigma_0^2}
 \le4r_t\frac{\sigma_w}{\sigma_0}.
 \end{aligned}
\]
\rbreak
Also, Cauchy--Schwarz gives
\(|\mathbb E[x_0w]|\le\sigma_0\sigma_w\), so
\[
 \begin{aligned}
 |\hat\beta_{0,\emptyset}-\beta_0|
 &=
 \frac{|x_0^\top w/n|}{x_0^\top x_0/n}\\
 &\le
 \frac{|\mathbb E[x_0w]|+r_t\sigma_0\sigma_w}
 {(1-r_t)\sigma_0^2}\\
 &\le
 \frac{1+r_t}{1-r_t}\frac{\sigma_w}{\sigma_0}
 \le3\frac{\sigma_w}{\sigma_0}.
 \end{aligned}
\]

For part (b), apply \eqref{eq:new-bilinear} with first vector
\((0,a^\top,0)^\top\), and second vector \(u\) when \(z=x_0\) or \(v\)
when \(z=w\).
\rbreak
Taking the supremum over \(a^\top C_{SS}a\le1\) gives
\begin{equation}\label{eq:new-vector-moments}
 \left\|
 C_{SS}^{-1/2}
 \left(\frac{X_S^\top z}{n}-\mathbb E[X_Sz]\right)
 \right\|
 \le r_t\|z\|_{L_2}.
\end{equation}

For the term involving the matrix inverse
\((X_S^\top X_S/n)^\dagger\), work on
\(\mathcal V_S=\{C_{SS}a:a\in\mathbb R^{|S|}\}\),
the span of the columns of \(C_{SS}\), with \(I\) the identity on
\(\mathcal V_S\), and set
\[
 G_S=C_{SS}^{-1/2}\frac{X_S^\top X_S}{n}C_{SS}^{-1/2}.
\]
\rbreak
Lemma~\ref{lem:new-uniform} gives
\(\|G_S-I\|_{\mathrm{op}}\le r_t\).
\rbreak
Hence, for every \(a\in\mathcal V_S\),
\begin{equation}\label{eq:new-inverse-norm}
 \begin{aligned}
 \|G_Sa\|
 &\ge\|a\|-\|(G_S-I)a\|
 \ge(1-r_t)\|a\|,\\
 \|G_S^{-1}\|_{\mathrm{op}}
 &\le\frac1{1-r_t}\le2.
 \end{aligned}
\end{equation}

The resolvent identity, applied with \(A=G_S\) and \(B=I\)
(\citeauthor{couilletLiao2022rmt}, Lemma~2.1), gives
\begin{equation}\label{eq:new-inverse-difference}
 \begin{aligned}
 G_S^{-1}-I&=G_S^{-1}(I-G_S),\\
 \|G_S^{-1}-I\|_{\mathrm{op}}
 &\le
 \|G_S^{-1}\|_{\mathrm{op}}\,
 \|I-G_S\|_{\mathrm{op}}
 \le\frac{r_t}{1-r_t}\le2r_t.
 \end{aligned}
\end{equation}
\rbreak
Further, on the same span \(\mathcal V_S\),
\[
 C_{SS}^{1/2}
 \left(\frac{X_S^\top X_S}{n}\right)^\dagger
 C_{SS}^{1/2}=G_S^{-1},\qquad
 C_{SS}^{1/2}C_{SS}^{-1}C_{SS}^{1/2}=I.
\]

Now write
\[
 \frac{x_0^\top P_Sz}{n}
 =
 \left(\frac{X_S^\top x_0}{n}\right)^\top
 \left(\frac{X_S^\top X_S}{n}\right)^\dagger
 \frac{X_S^\top z}{n}.
\]
\rbreak
Compare \(X_S^\top x_0/n\) with \(d_S\), the empirical inverse with
\(C_{SS}^{-1}\), and \(X_S^\top z/n\) with \(\mathbb E[X_Sz]\),
replacing one factor at a time.
\rbreak
Applying \(|a^\top B b|\le\|a\|\|B\|_{\mathrm{op}}\|b\|\)
to the resulting three differences gives the first inequality:
\begin{equation}\label{eq:new-projected-assembly}
 \begin{aligned}
 &\left|
 \frac{x_0^\top P_Sz}{n}
 -d_S^\top C_{SS}^{-1}\mathbb E[X_Sz]
 \right|\\
 &\le
 \left\|
 C_{SS}^{-1/2}\left(\frac{X_S^\top x_0}{n}-d_S\right)
 \right\|
 \|G_S^{-1}\|_{\mathrm{op}}
 \left\|C_{SS}^{-1/2}\frac{X_S^\top z}{n}\right\|\\
 &\quad+
 \|C_{SS}^{-1/2}d_S\|\,
 \|G_S^{-1}(I-G_S)\|_{\mathrm{op}}
 \left\|C_{SS}^{-1/2}\frac{X_S^\top z}{n}\right\|\\
 &\quad+
 \|C_{SS}^{-1/2}d_S\|
 \left\|
 C_{SS}^{-1/2}\left(\frac{X_S^\top z}{n}-\mathbb E[X_Sz]\right)
 \right\|\\
 &\le
 2r_t\left(\sigma_0+\sqrt{d_S^\top C_{SS}^{-1}d_S}\right)
 \left[
 \sqrt{\mathbb E[X_Sz]^\top C_{SS}^{-1}\mathbb E[X_Sz]}
 +r_t\|z\|_{L_2}
 \right]\\
 &\quad+
 r_t\|z\|_{L_2}\sqrt{d_S^\top C_{SS}^{-1}d_S}\\
 &\le
 4r_t\left[
 \sigma_0
 \sqrt{\mathbb E[X_Sz]^\top C_{SS}^{-1}\mathbb E[X_Sz]}
 +\|z\|_{L_2}\sqrt{d_S^\top C_{SS}^{-1}d_S}
 +r_t\sigma_0\|z\|_{L_2}
 \right],
 \end{aligned}
\end{equation}
where the second inequality uses
\eqref{eq:new-vector-moments}--\eqref{eq:new-inverse-difference}
and the triangle inequality, and the last uses
\(\sqrt{d_S^\top C_{SS}^{-1}d_S}\le\sigma_0\) and \(2r_t\le1\).
\end{proof}

Define
\begin{equation}\label{eq:new-psi}
 \Psi_s
 =
 \sqrt{
 \max_{\substack{J\subseteq[p]\\|J|\le s}}
 (C\beta_{-0})_J^\top C_{JJ}^{-1}(C\beta_{-0})_J
 }.
\end{equation}

\begin{theorem}\label{thm:new-search}
There exist a universal constant \(C_{\mathrm n}>0\) and a constant
\(C_{\mathrm u}>0\), depending only on \(K_X,K_\varepsilon\),
such that the following holds under Assumption~\ref{ass:new} and
\begin{equation}\label{eq:new-search-sample-size}
 n\ge
 C_{\mathrm n}(K_X^2+K_\varepsilon^2)^2
 s\log(ep/s).
\end{equation}

If \(d=0\), then, with probability at least
\(1-2\exp[-s\log(ep/s)]\),
\begin{equation}\label{eq:new-experimental}
 \max_{|S|\le s}
 |\hat\beta_{0,S}-\hat\beta_{0,\emptyset}|
 \le
 C_{\mathrm u}\frac{\delta_{s,n,p}}{\sigma_0}
 \left(\Psi_s+\delta_{s,n,p}\sigma_w\right).
\end{equation}

For general \(d\), let
\(S^\star\in\operatorname*{arg\,max}_{|S|=s}\beta^\star_{0,S}\).
With the same probability,
\begin{equation}\label{eq:new-observational}
 \begin{aligned}
 0
 &\le
 \max_{|S|=s}\hat\beta_{0,S}-\hat\beta_{0,S^\star}\\
 &\le
 C_{\mathrm u}\frac{\delta_{s,n,p}}{c_{\mathrm{ov}}\sigma_0}
 \left[
 \Psi_s+
 \sigma_w\left(
 \delta_{s,n,p}
 +\sqrt{
 \max_{|J|\le s}
 \frac{d_J^\top C_{JJ}^{-1}d_J}{\sigma_0^2}
 }
 \right)
 \right].
 \end{aligned}
\end{equation}
\end{theorem}

\begin{proof}
Work on \(\mathcal E_t\).
\rbreak
If \(d=0\), then \(b(S)=0\) and \(\eta(S)=\sigma_0^2\) for every \(S\).
\rbreak
Proposition~\ref{prop:new-coefficient} and \eqref{eq:new-psi} give
\begin{equation}\label{eq:new-experimental-radius}
 \max_{|S|\le s}
 |\hat\beta_{0,S}-\hat\beta_{0,\emptyset}|
 \le
 \frac{A r_t}{\sigma_0}
 \left(\Psi_s+r_t\sigma_w\right).
\end{equation}

For general \(d\), Assumption~\eqref{ass:new-overlap} and the definition
of \(\Psi_s\) give, for every \(|S|\le s\),
\[
 \begin{gathered}
 \eta(S)\ge c_{\mathrm{ov}}\sigma_0^2,\\
 \sqrt{(C\beta_{-0})_S^\top C_{SS}^{-1}(C\beta_{-0})_S}
 \le\Psi_s,\\
 \frac{\sqrt{d_S^\top C_{SS}^{-1}d_S}}{\sigma_0}
 \le
 \sqrt{
 \max_{|J|\le s}
 \frac{d_J^\top C_{JJ}^{-1}d_J}{\sigma_0^2}
 }.
 \end{gathered}
\]
\rbreak
Substituting these inequalities in
Proposition~\ref{prop:new-coefficient} yields, uniformly over \(|S|\le s\),
\begin{equation}\label{eq:new-observational-uniform}
 \begin{aligned}
 &\left|
 (\hat\beta_{0,S}-\hat\beta_{0,\emptyset})
 -\bigl(b(S)-b(\emptyset)\bigr)
 \right|\\
 &\quad\le
 \frac{A r_t}{c_{\mathrm{ov}}\sigma_0}
 \left[
 \Psi_s+
 \sigma_w\left(
 r_t+
 \sqrt{
 \max_{|J|\le s}
 \frac{d_J^\top C_{JJ}^{-1}d_J}{\sigma_0^2}
 }
 \right)
 \right].
 \end{aligned}
\end{equation}

By definition, \(S^\star\subseteq[p]\) and \(|S^\star|=s\), so
\[
 \max_{\substack{S\subseteq[p]\\|S|=s}}\hat\beta_{0,S}
 \ge\hat\beta_{0,S^\star}.
\]
\rbreak
Also, \(b(S)\le b(S^\star)\) whenever \(|S|=s\).
\rbreak
The triangle inequality and \eqref{eq:new-observational-uniform}
therefore imply
\begin{equation}\label{eq:new-observational-radius}
 \begin{aligned}
 0
 &\le
 \max_{|S|=s}\hat\beta_{0,S}-\hat\beta_{0,S^\star}\\
 &\le
 \max_{|S|=s}
 \left|
 (\hat\beta_{0,S}-\hat\beta_{0,\emptyset})
 -\bigl(b(S)-b(\emptyset)\bigr)
 \right|\\
 &\quad+
 \left|
 (\hat\beta_{0,S^\star}-\hat\beta_{0,\emptyset})
 -\bigl(b(S^\star)-b(\emptyset)\bigr)
 \right|\\
 &\le
 \frac{2A r_t}{c_{\mathrm{ov}}\sigma_0}
 \left[
 \Psi_s+
 \sigma_w\left(
 r_t+
 \sqrt{
 \max_{|J|\le s}
 \frac{d_J^\top C_{JJ}^{-1}d_J}{\sigma_0^2}
 }
 \right)
 \right].
 \end{aligned}
\end{equation}

Choose \(t=s\log(ep/s)\), which is at least one.
\rbreak
The dimension bound preceding \eqref{eq:new-radius-rate} gives
\[
 k+2+\log\binom pk+t
 \le7s\log(ep/s).
\]
\rbreak
Hence \eqref{eq:new-search-sample-size}, with
\(C_{\mathrm n}\ge7C_1\), implies the sample-size condition of
Lemma~\ref{lem:new-uniform}.
\rbreak
Moreover,
\begin{equation}\label{eq:new-radius-specialized}
 \begin{aligned}
 r_t^2
 &\le
 \frac{2C_0^2(K_X^2+K_\varepsilon^2)^2}{n}
 \left[k+2+\log\binom pk+t\right]\\
 &\le
 14C_0^2(K_X^2+K_\varepsilon^2)^2
 \delta_{s,n,p}^2.
 \end{aligned}
\end{equation}
\rbreak
Substituting this bound for \(r_t\) and \(r_t^2\) in
\eqref{eq:new-experimental-radius} and
\eqref{eq:new-observational-radius} proves
\eqref{eq:new-experimental} and \eqref{eq:new-observational},
with \(C_{\mathrm u}\) depending only on \(K_X,K_\varepsilon\).
\rbreak
Finally, Lemma~\ref{lem:new-uniform} gives
\[
 \mathbb P(\mathcal E_t)
 \ge1-2e^{-t}
 =
 1-2\exp[-s\log(ep/s)].\qedhere
\]
\end{proof}
\end{review}

\begin{thebibliography}{2}
\bibitem[Vershynin(2018)]{vershynin2018hdp}
Roman Vershynin.
\newblock \emph{High-Dimensional Probability: An Introduction with
Applications in Data Science}.
\newblock Cambridge University Press, 2018.
\newblock First edition, Theorem~4.6.1 and equation~(4.22).
\newblock \href{https://webapps.math.uci.edu/~rvershyn/papers/HDP-book/HDP-1.pdf}{Author's online edition}.
\bibitem[Couillet and Liao(2022)]{couilletLiao2022rmt}
Romain Couillet and Zhenyu Liao.
\newblock \emph{Random Matrix Methods for Machine Learning}.
\newblock Cambridge University Press, 2022.
\newblock Lemma~2.1.
\newblock \href{https://polaris.imag.fr/romain.couillet/docs/RMT_ML_Book.pdf}{Authors' online edition}.
\end{thebibliography}

\end{document}
